\chapter{Marco Teórico}

\section{Sistemas de punto de venta}

Un sistema de punto de venta (POS, por sus siglas en inglés) registra la operación comercial en el momento en que se realiza. El ciclo básico comprende identificar artículos, determinar cantidades y precios, calcular el importe, registrar el medio de pago y conservar el comprobante de la transacción. En Dulcería POS, el servidor valida de nuevo los productos y sus precios antes de guardar la venta; de esta forma el importe persistido no depende exclusivamente de lo que envíe el navegador.

El registro relaciona una cabecera de venta con varios renglones de detalle. La cabecera conserva folio, usuario, fecha, método de pago e importes; cada renglón conserva producto, cantidad, precio unitario e importe. Esa estructura permite reconstruir el contenido de un ticket y analizar ventas por periodo.

\section{Gestión de inventario y caja}

El inventario representa la cantidad disponible de cada producto. Para una venta con cantidad \(q_i\) del producto \(i\), el saldo esperado tras la operación es:
\[
  S_{i,\mathrm{nuevo}} = S_{i,\mathrm{anterior}} - q_i,
\]
donde \(S_i\) es el stock registrado. El sistema compara las existencias con el umbral mínimo configurado y considera bajo el stock cuando \(S_i \leq S_{i,\min}\). El código permite registrar stock decimal en el modelo; la validación del POS web restringe actualmente la cantidad vendida a números enteros.

El control de caja distingue efectivo y tarjeta. El efectivo esperado del periodo considera las ventas en efectivo, los ingresos y retiros manuales y el fondo de caja dejado por el corte anterior:
\[
  E_{\mathrm{esperado}} =
  V_{\mathrm{efectivo}} + I_{\mathrm{caja}} - R_{\mathrm{caja}} + F_{\mathrm{anterior}}.
\]
El total por tarjeta se calcula por separado a partir de las ventas correspondientes. Si se proporciona el efectivo contado, la diferencia se obtiene mediante:
\[
  D_{\mathrm{caja}} = E_{\mathrm{contado}} - E_{\mathrm{esperado}}.
\]
Un valor positivo indica sobrante y uno negativo indica faltante. El cálculo se centraliza en el servicio \texttt{CajaResumen}, reutilizado por las pantallas de caja, cortes y dashboard.

\section{Arquitectura web y separación por empresa}

La arquitectura cliente-servidor separa la presentación en navegador de las reglas y persistencia del servidor. El frontend de Vue 3 se comunica con Laravel mediante Inertia.js para las pantallas web; Laravel valida solicitudes, aplica permisos y coordina las operaciones sobre MySQL. La API sigue rutas REST y protege las solicitudes autenticadas con tokens de Laravel Sanctum. La documentación oficial de Laravel describe sus mecanismos de rutas, middleware, autenticación y persistencia \cite{laravel12,sanctum}; Vue documenta el modelo de componentes reactivos utilizado por la interfaz \cite{vue3}.

La estrategia multiempresa asigna un identificador \texttt{tenant\_id} a los usuarios y registros de negocio. El trait \texttt{HasTenant} aplica un filtro global a los modelos que lo utilizan y completa el tenant al crear registros con un usuario autenticado. Las rutas restringidas y los middleware aplican además las reglas de rol. La separación lógica de datos no equivale por sí misma a una base de datos física independiente: el diseño actual usa tablas compartidas con registros asociados a cada tenant.

\section{Persistencia relacional y transacciones}

Una base de datos relacional organiza la información en tablas vinculadas por claves primarias y foráneas. En este proyecto, una venta pertenece a un usuario y puede tener varios detalles; cada detalle apunta a un producto. El inventario, los usuarios, los tenants y los movimientos de caja se relacionan para conservar contexto y permitir consultas.

El registro de una venta agrupa sus escrituras en una transacción. Dentro de ella, los productos se vuelven a consultar con bloqueo de fila, se comprueba que estén activos y que tengan existencias, se calcula el precio desde la base de datos y se comprueba que el total coincida con el que confirmó el cajero. A continuación, se crea la venta y sus detalles, y se decrementa el stock. Una transacción permite confirmar el conjunto como una operación indivisible ante errores de persistencia \cite{mysqltransactions}.

\section{Control de acceso y trazabilidad}

La autenticación comprueba la identidad; la autorización determina qué acciones puede realizar esa identidad. Dulcería POS implementa comprobaciones de usuario activo, middleware de administración y superadministración, y restricciones adicionales para que un cajero consulte únicamente sus ventas y cortes. La API utiliza tokens de acceso de Sanctum y límites de solicitudes en rutas protegidas. El sistema registra eventos de auditoría para acciones seleccionadas, incluyendo inicio y cierre de sesión.

\section{Proceso de desarrollo de software}

El desarrollo incremental organiza el producto en módulos que pueden construirse e integrarse progresivamente. Sommerville presenta los procesos de software como actividades estructuradas de especificación, desarrollo, validación y evolución \cite{sommerville2015}. Para este proyecto, este enfoque facilita priorizar primero el flujo crítico de venta y después incorporar administración, caja, reportes y supervisión, verificando la integración con las rutas y modelos existentes.

\section{Tecnologías utilizadas}

\begin{table}[htbp]
\centering
\caption{Tecnologías que integran la solución}
\label{tab:tecnologias}
\begin{tabularx}{\textwidth}{|>{\raggedright\arraybackslash}p{3.1cm}|>{\raggedright\arraybackslash}p{3.4cm}|X|}
\hline
\textbf{Componente} & \textbf{Tecnología} & \textbf{Uso en el proyecto} \\
\hline
Servidor y lógica & PHP 8.2 o superior, Laravel 12 & Rutas web y API, validación, autorización, modelos y transacciones. \\
\hline
Interfaz web & Vue 3, Inertia.js & Vistas reactivas servidas desde la aplicación Laravel. \\
\hline
Datos & MySQL 8 & Persistencia relacional de tenants y operaciones. \\
\hline
Autenticación API & Laravel Sanctum & Tokens Bearer para las rutas REST protegidas. \\
\hline
Entrega & Docker Compose, Nginx y PHP-FPM & Ejecución reproducible de la aplicación y sus servicios. \\
\hline
\end{tabularx}
\end{table}